\chapter[Planejamento do Experimento]{Planejamento do Experimento}

Este capítulo apresenta a proposta de planejamento do experimento controlado que será realizado para avaliar o impacto de 
diferentes padrões de engenharia de prompt na aferição do indicador de manutenibilidade, considerando modelos de 
qualidade de referência de produto de software, no contexto de projetos de código-fonte aberto para aplicações web. 

Ao longo do capítulo, são descritos os principais elementos que compõem a estrutura experimental adotada. Inicialmente, 
apresenta-se a definição do experimento, delimitando seu propósito, escopo e relação com o problema de pesquisa. Em seguida, 
é realizada a seleção do contexto, com a caracterização do ambiente experimental, dos projetos avaliados e das condições sob as 
quais a investigação será conduzida.

Também são formuladas as hipóteses que orientarão a análise dos resultados, bem como as variáveis independentes, dependentes e 
de controle consideradas no estudo. Na sequência, são apresentados a seleção dos sujeitos e a definição dos objetos experimentais, 
compreendendo os artefatos de software que serão utilizados na avaliação.

O capítulo ainda descreve o design do experimento, incluindo a organização dos tratamentos, a estratégia de execução e os procedimentos
previstos para coleta dos dados. Posteriormente, são definidos os critérios de avaliação, com foco na forma como os resultados serão  
analisados em relação à manutenibilidade e aos modelos de qualidade adotados como referência. Por fim, são apresentados os próximos passos do experimento, 
indicando as atividades necessárias para sua execução, validação e posterior análise dos resultados.

\section{Definição}

A abordagem experimental proposta para esta pesquisa baseia-se em um experimento controlado em Engenharia de Software. A escolha por essa abordagem decorre da 
à partir da estrutura da questão de pesquisa (verificar em \ref{sec:questao_de_pesquisa}). Nesse sentido, a pesquisa não se limita à descrição de um fenômeno, mas propõe a comparação entre diferentes 
tratamentos aplicados a um mesmo tipo de objeto experimental, com o objetivo de observar seus efeitos sobre uma variável de resultado.

Segundo \citeonline{wohlin2024}, experimentos são instrumentos relevantes para avaliar métodos, técnicas, linguagens, ferramentas ou soluções em Engenharia de Software, 
especialmente quando se deseja comparar alternativas e obter dados quantitativos que apoiem uma análise empírica. Os autores destacam que a experimentação contribui para 
tornar a Engenharia de Software uma disciplina mais científica, pois permite avaliar propostas de pesquisa a partir de hipóteses, evidências e observações sistemáticas. 
Essa perspectiva está alinhada ao propósito deste trabalho, uma vez que se pretende avaliar, de forma controlada, se diferentes estratégias de prompt produzem resultados 
distintos na avaliação da manutenibilidade.

\citeonline[p.~73, tradução nossa]{wohlin2024} definem um experimento controlado em Engenharia de Software como uma investigação empírica que manipula "um fator ou variável do cenário estudado". Ainda segundo os autores, diferentes tratamentos são aplicados aos sujeitos ou objetos, enquanto outras variáveis são mantidas constantes, permitindo 
medir os efeitos produzidos sobre as variáveis de resultado \cite[p.~73]{wohlin2024}. Essa definição é diretamente aplicável ao desenho deste trabalho, pois os padrões de engenharia 
de prompt representam o fator manipulado, enquanto os resultados de aferição da manutenibilidade representam a variável dependente observada.

A escolha pelo experimento controlado também se justifica porque o estudo busca avaliar uma possível relação entre tratamento e resultado. De acordo com \citeonline{wohlin2024}, um experimento 
parte de uma hipótese ou de uma relação esperada entre causa e efeito; a partir disso, os tratamentos são aplicados, os resultados são observados e, caso o experimento seja adequadamente planejado, 
torna-se possível analisar a relação entre a causa manipulada e o efeito observado. No contexto desta pesquisa, a causa investigada é a variação dos padrões de prompt, enquanto o efeito observado 
é a variação na aferição da manutenibilidade produzida pelos LLMs.

Portanto, a adoção de um experimento controlado é coerente com o desenho desta pesquisa, pois permite manipular a variável de interesse, controlar fatores externos, comparar tratamentos e analisar 
os efeitos observados sobre uma variável dependente. Essa abordagem contribui para que o trabalho prático seja conduzido de forma sistemática, transparente e passível de replicação, fortalecendo 
a validade da investigação e a confiabilidade dos resultados obtidos.

\section{Seleção de Contexto}

A seleção do contexto define o ambiente no qual o experimento será conduzido e delimita as condições sob as quais os resultados deverão ser interpretados. De acordo com \citeonline[p.~110, tradução nossa]{wohlin2024}, 
o contexto corresponde ao "ambiente no qual o experimento é executado", contemplando tanto os sujeitos envolvidos quanto os artefatos de software utilizados como objetos do experimento.

No presente trabalho, o estudo será conduzido no contexto da avaliação de qualidade de produto de software com apoio de Modelos de Linguagem de Grande Escala, com foco específico na característica de manutenibilidade. 
A escolha desse contexto, portanto, está alinhada à questão de pesquisa (verificar em \ref{sec:questao_de_pesquisa}).

Dessa forma, o experimento será realizado em um ambiente controlado, utilizando os mesmos objetos de análise para diferentes variações de prompt, de modo a permitir comparação entre os resultados produzidos 
com modelos de referência de avaliação de qualidade de software.

\section{Formulação de Hipóteses}

A formulação das hipóteses parte da premissa de que diferentes padrões de engenharia de prompt podem influenciar a forma como um Modelo de Linguagem de Grande Escala avalia a manutenibilidade de trechos de código-fonte. 
Como o experimento busca comparar resultados obtidos a partir de diferentes tratamentos, as hipóteses foram definidas de modo a verificar se a variação dos padrões de prompt produz diferenças significativas nos resultados de avaliação.

Segundo \citeonline{wohlin2024}, a análise estatística de um experimento é fundamentada no teste de hipóteses. Os autores explicam que uma hipótese deve ser formalmente estabelecida e que os dados coletados durante o experimento 
são utilizados para verificar se há evidências suficientes para rejeitá-la. Ainda segundo os autores, a hipótese nula representa a ausência de tendências ou padrões reais no ambiente experimental, de modo que eventuais diferenças 
observadas seriam atribuídas ao acaso \cite{wohlin2024}.

Neste trabalho, seja $M$ a medida utilizada para representar o resultado da avaliação de manutenibilidade produzida pelo LLM, e sejam $P_1, P_2, \ldots, P_n$ os diferentes padrões de engenharia de prompt avaliados. A hipótese 
nula considera que os padrões de prompt apresentam desempenho equivalente na avaliação de manutenibilidade, isto é, não há diferença significativa entre os resultados produzidos pelos diferentes padrões. A hipótese alternativa 
considera que existe diferença significativa entre os resultados produzidos por pelo menos dois padrões de prompt.

Dessa forma, as hipóteses do experimento são formuladas da seguinte maneira:

\begin{itemize}
    \item \textbf{H0}: diferentes padrões de engenharia de prompt não produzem diferenças significativas na avaliação de manutenibilidade realizada por LLMs.
    
    \[
    H_0: \mu_M(P_1) = \mu_M(P_2) = \ldots = \mu_M(P_n)
    \]

    \item \textbf{H1}: diferentes padrões de engenharia de prompt produzem diferenças significativas na avaliação de manutenibilidade realizada por LLMs.
    
    \[
    H_1: \exists\, P_i, P_j \mid \mu_M(P_i) \neq \mu_M(P_j)
    \]
\end{itemize}

Assim, a hipótese nula afirma que a média dos resultados de manutenibilidade é equivalente entre os padrões de prompt avaliados. Já a hipótese alternativa afirma 
que ao menos um par de padrões de prompt apresenta diferença significativa em relação à medida de manutenibilidade observada.

A avaliação dessas hipóteses será realizada a partir dos resultados gerados pelos LLMs para os mesmos objetos experimentais, submetidos a diferentes padrões de prompt. 
Dessa forma, busca-se verificar se as variações observadas nos resultados podem ser associadas aos tratamentos aplicados ou se permanecem dentro de uma variação 
esperada entre avaliações equivalentes.

\section{Seleção de Variáveis}

\section{Seleção dos Sujeitos}

\subsection{Definição dos Objetos}

\section{Design do Experimento}

\section{Avaliação}

\section{Próximos passos do Experimento}